En la industria automotriz y el automovilismo deportivo modernos, el prototipado digital de vehículos ha adquirido una importancia fundamental. El constante aumento en los costos de desarrollo y fabricación, sumado a la creciente complejidad de los sistemas vehiculares, exige soluciones eficientes. Por ello, la industria está migrando aceleradamente hacia el uso de simulaciones avanzadas. Estas herramientas permiten evaluar diferentes enfoques de diseño en las etapas tempranas de construcción, lo que reduce drásticamente los costos, facilita la detección prematura de errores y permite analizar con precisión cómo se comporta dinámicamente el sistema en su conjunto.\\

Además de su valor en el diseño, este tipo de simulaciones se utiliza extensamente como herramienta de entrenamiento y preparación para pilotos de carreras (\textit{Driver-in-the-Loop}). Estos sistemas les permiten utilizar un entorno virtual que replica fielmente las condiciones reales de pista, a una fracción del costo y eliminando los riesgos físicos y materiales que implicaría realizar pruebas en un circuito.\\

Aunque existen plataformas comerciales robustas para llevar a cabo estas simulaciones, actualmente se observa una falta de alternativas de código abierto que integren las funcionalidades requeridas empleando técnicas y lenguajes de modelado modernos. Las soluciones en el mercado suelen ser propietarias y cerradas, lo que representa una barrera de entrada y limita severamente su adaptación, extensión o auditoría por parte de investigadores y desarrolladores independientes.\\

A esta limitación se suma la necesidad crítica de contar con una arquitectura de software inherentemente modular. El estudio de la dinámica vehicular exige una plataforma flexible, capaz de escalar en complejidad para satisfacer distintos objetivos de análisis o adaptarse a diferentes capacidades de cómputo. Para resolver este desafío, resulta coherente la adopción de Modelica. Gracias a su paradigma de modelado orientado a objetos, este lenguaje permite la encapsulación de cada componente físico, facilitando el intercambio fluido de los distintos subsistemas del vehículo y sentando las bases para una herramienta verdaderamente extensible e interoperable.\\
